\documentclass[10pt,a4paper]{article}
\usepackage{amsmath,amssymb,amsthm,graphicx}
\usepackage{amsmath,amsthm,amssymb,mathtools,graphicx,xcolor,textcomp,multirow}
\usepackage[no-math]{fontspec}
\usepackage[linesnumbered,ruled,procnumbered]{algorithm2e}
\usepackage[colorlinks=true,hypertexnames=false]{hyperref}
\usepackage[capitalize]{cleveref}
\newtheorem{theorem}{Theorem}
\newtheorem{lemma}[theorem]{Lemma}
\newtheorem{definition}[theorem]{Definition}
\newtheorem{notation}[theorem]{Notation}
\newtheorem{proposition}[theorem]{Proposition}
\newtheorem{corollary}[theorem]{Corollary}
\newtheorem{conjecture}[theorem]{Conjecture}
\newtheorem{assumption}[theorem]{Assumption}
\newtheorem{fact}[theorem]{Fact}
\newtheorem{remark}[theorem]{Remark}
\newtheorem{claim}[theorem]{Claim}
\newtheorem{example}[theorem]{Example}
\newtheorem{problem}[theorem]{Problem}
\newtheorem{hypothesis}[theorem]{Hypothesis}
\newtheorem{observation}[theorem]{Observation}
\newcommand{\wh}{\widehat}
\newcommand{\wt}{\widetilde}
\newcommand{\ov}{\overline}
\newcommand{\eps}{\epsilon}
\newcommand{\N}{\mathbb{N}}
\newcommand{\A}{\mathcal{A}}

\newcommand{\Hil}{\mathcal{H}}
\newcommand{\F}{\mathbb{F}}
\newcommand{\R}{\mathbb{R}}
\renewcommand{\P}{\mathbb{P}}
\newcommand{\RHS}{\mathrm{RHS}}
\newcommand{\LHS}{\mathrm{LHS}}
\renewcommand{\d}{\mathrm{d}}
\renewcommand{\i}{\mathbf{i}}
\renewcommand{\varepsilon}{\lambda}
\renewcommand{\tilde}{\wt}
\renewcommand{\hat}{\wh}
\renewcommand{\eps}{\epsilon}
\newcommand{\sig}{\mathrm{sig}}
\newcommand{\high}{\mathrm{high}}
\newcommand{\bad}{\mathrm{bad}}
\newcommand{\diag}{\textrm{diag}}
\renewcommand{\d}{\mathrm{d}}
\newcommand{\ip}[2]{\langle {#1} , {#2} \rangle}
\newcommand{\norm}[1]{\| #1 \|}
\DeclareMathOperator*{\E}{{\mathbb{E}}}
\DeclareMathOperator*{\Var}{{\bf {Var}}}
\DeclareMathOperator*{\Sup}{{\bf {Sup}}}
\DeclareMathOperator*{\Z}{\mathbb{Z}}
%
\DeclareMathOperator*{\C}{\mathbb{C}}
\DeclareMathOperator{\supp}{supp}
\newcommand{\epsf}{\eta}
\newcommand{\ket}[1]{\left| #1 \right\rangle}
\newcommand{\bra}[1]{\left\langle #1 \right|}

\DeclareMathOperator{\poly}{poly}
\DeclareMathOperator{\polylog}{polylog}
\DeclareMathOperator{\nnz}{nnz}
\DeclareMathOperator{\rank}{rank}
\DeclareMathOperator{\new}{new}
\DeclareMathOperator{\pre}{pre}
\DeclareMathOperator{\old}{old}
\DeclareMathOperator{\tr}{tr}
\DeclareMathOperator*{\argmin}{arg\,min}
\newcommand{\ceil}[1]{\lceil #1 \rceil}
\newcommand{\floor}[1]{\lfloor #1 \rfloor}

\renewcommand{\S}{\mathcal{S}}
\newcommand{\U}{\mathcal{U}}
\newcommand{\id}{\mathrm{id}}



%
\newcommand{\ora}{\mathcal{O}}
\newcommand{\FG}{\mathrm{FG}}
\newcommand{\QFG}{\mathrm{QFG}}
\newcommand{\CP}{\ensuremath{\mathsf{CP}}}
\newcommand{\OV}{\ensuremath{\mathsf{OV}}}
\newcommand{\OVH}{\ensuremath{\mathsf{OVH}}}
\newcommand{\SAT}{\ensuremath{\mathsf{SAT}}}
\newcommand{\BCP}{\ensuremath{\mathsf{BCP}}}
\newcommand{\xiBCP}{(1+\xi)\text{-}\mathsf{BCP}}
\newcommand{\xiCP}{(1+\xi)\text{-}\mathsf{CP}}
\newcommand{\eCP}{\mathsf{CP}_{\epsilon}}
\newcommand{\eBCP}{\mathsf{BCP}_{\epsilon}}
\newcommand{\xieBCP}{(1+\xi)\text{-}\mathsf{BCP}_{\epsilon}}
\newcommand{\ZOV}{\ensuremath{\Z\text{-}\mathsf{OV}}}
\newcommand{\CRR}{\mathsf{CRR}}
\newcommand{\MA}{\mathsf{MA}}

\crefname{claim}{Claim}{Claims}
\crefname{procedure}{Procedure}{Procedures}


%
\SetAlgoProcName{Procedure}{Procedure}


\newcommand{\zo}{\{0,1\}}
\SetAlgoProcName{Procedure}{Procedure}
\allowdisplaybreaks
\AtBeginDocument{\author{\parbox{.95\textwidth}{\centering Scott Aaronson, Nai-Hui Chia, Han-Hsuan Lin,\\ Chunhao Wang, Ruizhe Zhang}}}
\title{QSETH-conditional quantum lower bounds for high-dimensional closest pair and related geometric search}
\author{Scott Aaronson, Nai-Hui Chia, Han-Hsuan Lin, Chunhao Wang, Ruizhe Zhang}
\date{doi:10.4230/LIPIcs.CCC.2020.16; CC BY 3.0 Unported}
\begin{document}
\maketitle
\noindent\textit{Editorial scope.} One conditional quantum fine-grained hardness mechanism with full SAT/OV/BCP/CP and related integer-OV reduction family. Related dimensions/exponents and bichromatic consequences grouped once. QSETH is assumed;this task does not assert an unconditional separation or duplicate constant-dimensional algorithm GA101.\par\medskip
\noindent\textit{Allowed standard prerequisites (editorial).}\begin{itemize}
\item QSETH,Grover/minimum finding,known contact-dimension realization/covering and integer-OV encodings. QSETH remains an explicit assumption,not a theorem. Formal Theorem31/Lemma33[29] and Theorem37[17] retained as cited prerequisite constructions with full stated hypotheses;all author's quantum-oracle realization,cost/probability analysis and related reductions are retained. Full52 original bibliography entries retained.
\end{itemize}\medskip
\setcounter{section}{1}
\setcounter{theorem}{5}
\section{Preliminaries}
\begin{definition}[Distance measure]
For any two vectors $a,b\in\R^d$, the distance between them in the $\ell_2$-metric is denoted by $\|a-b\|=\left(\sum_{i=1}^d|a_i-b_i|^2\right)^{1/2}$.
Their distance in the $\ell_0$-metric (Hamming distance) is denoted by 
$\|a-b\|_0=|\{i\in[d]:a_i\neq b_i\}|$, i.e.,
the number of coordinates on which $a$ and $b$ differ.
\end{definition}


\subsection{Quantum query model}

We consider the quantum query model in this work. Let $X:=\{x_1,\dots,x_n\}$ be a set of $n$ input points and $\ora_X$ be the corresponding oracle. We can access the $i$-th data point $x_i$ by making the query 
\begin{align}
    \ket{i}\ket{0} \xrightarrow{\ora_X} \ket{i}\ket{x_i},  
\end{align}
and we can make queries to elements in $X$ in superposition. Note that $\ora_X$ is an unitary transformation in the formula above. Hence, a quantum algorithm with access to $\ora_X$ can be represented as a sequence of unitary transformations. 

Consider a quantum algorithm $\A$ with access to an oracle $\ora$ and a initial state $\ket{0}:= \ket{0}_{Q}\ket{0}_{A}\ket{0}_W$, where the registers $Q$ and $A$ are for the queries and the answers from the oracle, and the register $W$ is the working space which is always hold by $\A$.  Then, we can represent the algorithm as 
\begin{align}
    U_T \ora U_{T-1}\cdots\ora U_1  \ket{0}. 
\end{align}
Let $\ket{\psi}_i= U_i\ora \cdots \ora U_1\ket{0}:= \sum_{i,z} \ket{i}_Q\ket{0}_A\ket{z}_W$ be the state right before applying the $i$-th $\ora$, then
\begin{align}
    \ora \ket{\psi}_i:= \sum_{i,z} \ket{i}_Q\ket{x_i}_A\ket{z}_W.
\end{align}
\subsection{Quantum subroutine for unstructured searching and minimum finding}

\begin{definition}[Unstructured search]
Given a set $P$ of $n$ elements in $\{0,1\}$, decide whether there exists a $1$ in $P$. 
\end{definition}

\begin{theorem}[Grover's search algorithm~\cite{r24,r37}]\label{thm:grover_algo}
There is a quantum algorithm for unstructured search with running time $O(\sqrt{n})$.
\end{theorem}

By \cref{thm:grover_algo} and BBBV's argument \cite{r9}, the quantum time complexity of unstructured search is $\Theta(\sqrt{n})$. We can also get a $\tilde{O}(\sqrt{n})$ quantum algorithm for minimum finding by combining Grover's search algorithm and binary search.
\begin{theorem}[Quantum minimum finding~\cite{r23}]\label{thm:mim}
There is a quantum algorithm that finds from a set of $n$ elements with values in $\mathbb{R}$, the index of the minimum element of the set, with success probability $\frac{1}{2}$ and run time $\tilde{O}(\sqrt{n})$.
\end{theorem}

\subsection{Problem definitions}
In this subsection, we first formally define  $\OV$, $\CP$, and $\BCP$. Then we show the folklore algorithms for $\CP$, $\BCP$, and $\OV$ by Grover's algorithm, which run in time $\tilde{O}(n)$. 
\begin{definition}[Orthogonal Vectors, $\OV$]\label{def:OV}
Given two sets $A,B$ of $n$ vectors in $\mathbb \{0,1\}^d$ as input, find a pair of vectors $a\in A$, $b\in B$ such that
$\ip{a}{b}=0$, where the inner product is taken in $\mathbb{Z}$.\footnote{Our definition is slightly different than some of the literature, for example, \cite{r18}, which is searching among pairs inside one set. Those two definitions are equivalent up to constant in complexities.}
\end{definition}

We denote $\OV$ with input length $n$ and dimension $d$ as $\OV_{n,d}$. We will use this notation when we need to specify the parameters in the following sections.  

\begin{definition}[Closest Pair Problem, $\CP$]
Given a set $P$ of $n$ points in $\mathbb R^d$ and a distance measure $\Delta$, find a pair of distinct points $a,b\in P$ such that $\Delta(a,b)$ is the smallest among all distinct pairs in $P$.
\end{definition}
Similar to $\OV$, we denote $\CP$ with input length $n$ and dimension $d$ as $\CP_{n,d}$. We will use this notation when the parameters in the following sections are required to be specified. Note that in this work, we consider $\Delta(a, b) = \norm{a-b}$ as the distance measure for \textsf{CP} and \textsf{BCP}.

\begin{definition}[Bichromatic Closest Pair Problem, \textsf{BCP}]
Given two sets $A,B$ of $n$ points in $\mathbb R^d$ and a distance measure $\Delta$, find a pair of points $a\in A$, $b\in B$ such that 
\begin{align}
    \Delta(a,b)= \min_{a\in A,b\in B} \Delta(a,b).
\end{align}
\end{definition}
We also define an approximate version of $\BCP$ as follows.
\begin{definition}[$(1+\xi$)-approximate Bichromatic Closest Pair Problem, $\xiBCP$]
Given two sets $A,B$ of $n$ points $\in \mathbb R^d$ and a distance measure $\Delta$, find a pair of points $a\in A$, $b\in B$ such that 
\begin{align}
    \Delta(a,b) \leq (1+\xi)\min_{a\in A,b\in B} \Delta(a,b).
\end{align}
\end{definition}
Same as $\CP$, we use $\BCP_{n,d}$ and $\xiBCP_{n, d}$ to specify the parameters. 

\begin{definition}[Element Distinctness Problem, \textsf{ED}]
  Let $f:[n] \rightarrow [m]$ be a given function. Decide whether there exist distinct $i, j\in[n]$ such that $f(i) = f(j)$.
\end{definition}

For this problem, Ambainis \cite{r5} gave a quantum algorithm with time complexity $\tilde{O}(n^{2/3})$, which matches the lower bound proved by Aaronson and Shi \cite{r1} up to a polylogarithmic factor.

\begin{theorem}\label{cor:trivial_alg}
There are $\tilde{O}(n)$-time quantum algorithms for $\CP$ and $\BCP$ when $d=O(\poly\log n)$.
\end{theorem}
\begin{proof}
We can solve $\CP$ and $\BCP$ by searching the minimum distance through all pairs by the algorithm of \cref{thm:mim}. There are $O(n^2)$ pairs and checking each pair took $O(d)$ time, so the total running time is $O(nd)$. For $d=O(\poly\log n)$, the time complexity equals to $\tilde{O}(n)$.
\end{proof}

\subsection{Fine-grained complexity}

As we have mentioned earlier in the introduction, a fine-grained reduction from problem $\mathsf{P}$ to $\mathsf{Q}$ with conjectured lower bounds $p(n)$ and $q(n)$, respectively, has the property that if we can improve the $q(n)$ time for $\mathsf{Q}$, then we can also improve the $p(n)$ time for $\mathsf{P}$. We give the formal definition by Williams~\cite{r46} in below. 

\begin{definition}[Fine-grained reduction, \cite{r46}]\label{def:fg}
Let $p(n)$ and $q(n)$ be non-decreasing functions of $n$. Problem $\mathsf{P}$ is ($p,q$)-reducible to problem $\mathsf{Q}$, denoted as $(\mathsf{P},p)\leq_{\FG} (\mathsf{Q},q)$, if for every $\epsilon$, there exist $\delta>0$, an algorithm $R$ for solving $\mathsf{P}$ with access to an oracle for $\mathsf{Q}$, a constant $d$, and an integer $k(n)$, such that for every $n\ge 1$, the algorithm $R$ takes any instance of $\mathsf{P}$ of size $n$ and 
\begin{itemize}
    \item $R$ runs in at most $d\cdot (p(n))^{1-\delta}$-time, 
    \item $R$ produces at most $k(n)$ instances of $\mathsf{Q}$ adaptively, that is, the $j$th instance $X_j$ is a function of $\{(X_i,y_i)\}_{1\leq i<j}$ where $X_i$ is the $i$th instance produced and $y_i$ is the answer of the oracle for $\mathsf{Q}$ on instance $X_i$, and 
    \item the sizes $n_i$ of the instances $X_i$ for any choice of oracle answers $y_i$ obeys the inequality
    \begin{align}\label{eq:fgr}
        \sum_{i=1}^{k(n)} (q(n_i))^{1-\epsilon} \leq d\cdot (p(n))^{1-\delta}.
    \end{align}
\end{itemize}
\end{definition}

Let $(\mathsf{P},p)\leq_{\FG} (\mathsf{Q},q)$ for some non-decreasing function $p(n)$ and $q(n)$. If for every $\eps > 0$, we can solve problem $\mathsf{Q}$ in time $q(n)^{1-\epsilon}$ with probability $1$ for all input length $n$, then there exists a $\delta>0$ such that we can solve the problem $\mathsf{P}$ in time $p(n)^{1-\delta}$ 
by \cref{eq:fgr}.

Here are some known results about fine-grained reductions.
\begin{theorem}[\cite{r29,r49}]\label{thm:classical_fg}
\begin{align}
    (\mathsf{CNF\text{-}SAT}_{n},2^n)\leq_{\FG} (\OV_{n_1,d_1},n_1^2)\leq_{\FG} (\BCP_{n_2,d_2},n_2^2)\leq_{\FG} (\CP_{n_3,d_3},n_3^2),
\end{align}
where $d_1=\Theta(\log n_1)$, $d_2=\Theta(\log n_2)$ and $d_3=(\log  n_3)^{\Omega(1)}$.
\end{theorem}

\begin{remark}
The second reduction from $\OV$ to $\BCP$ has been improved to $d_2 = 2^{O(\log^* n)}$ by Chen~\cite{r17}.
\end{remark}

There are several plausible hypotheses in fine-grained complexity, which can imply conditional hardness results for many interesting problems. We first give the definition of the strong exponential time hypothesis (SETH). 
\begin{hypothesis}[Strong Exponential Time Hypothesis, \textsf{SETH}]
For every $\epsilon > 0$, there exists a $k = k(\epsilon) \in \N$ such that no algorithm can solve $k$-$\SAT$ (i.e., satisfiability on a CNF of width $k$) in $O(2^{(1 - \epsilon)m})$ time where $m$ is the number of variables. Moreover, this holds even when the number of clauses is at most $c(\epsilon) \cdot m$ where $c(\epsilon)$ denotes a constant that depends only on $\epsilon$.
\end{hypothesis}

Another popular conjecture is the orthogonal vector hypothesis (\textsf{OVH}):
\begin{definition}[Orthogonal Vector Hypothesis, \textsf{OVH}]
For every $\epsilon > 0$,  
there exists a $c\ge 1$ such that $\OV_{n,d}$ requires $n^{2-\epsilon}$ time when $d=c\log n$.
\end{definition}

\begin{remark}
Under SETH, we can have the following conclusions from \cref{thm:classical_fg}: 
\begin{itemize}
    \item $\OVH$ is true. 
    \item For all $\epsilon>0$, there exists a $c>0$ such that $\BCP_{n,c\log n}$ cannot be solved by any randomized algorithm in time $O(n^{2-\epsilon})$.
    \item For all $\epsilon>0$, there exists a $c>0$ such that $\CP_{n,(\log n)^c}$ cannot be solved by any randomized algorithm in time $O(n^{2-\epsilon})$.
\end{itemize}
\end{remark}

\subsection{The framework for quantum walk search}
\label{sec:qw-framework}
In this subsection, we review the quantum walk framework for the Markov chain search problem and demonstrate how to use it to solve the element distinctness problem. For simplicity, we use the transition matrix $P$ to refer to a Markov chain, where $P = (p_{xy})_{x, y \in X}$ for $X$ being the state space of $P$ and $p_{xy}$ being the transition probability from $x$ to $y$. An irreducible and ergodic Markov chain has a unique stationary distribution $\pi$, which is also the unique eigenvector of $P$ with eigenvalue 1. Let $M \subseteq X$ be a set of marked elements. In the Markov chain search problem, the objective is to find an $x \in M$. We can perform the following actions: setup, sampling from the $\pi$ with cost $\mathsf{S}$; update, making a transition with cost $\mathsf{U}$, and checking whether the current state is marked or not with cost $\mathsf{C}$. To solve the search problem classically, we perform a random walk as follows. Start from a point sampled from $\pi$ and check if it is marked. If not, make a number of transitions on $P$ until it mixes, and then check again. We then repeat this process until a marked state is found. The cost of this random walk algorithm is $O(\mathsf{S} + \frac{1}{\varepsilon}(\frac{1}{\delta}\mathsf{U} + \mathsf{C}))$, where $\varepsilon:= |M|/|X|$ and $\delta$ is the spectral gap of $P$.

Quantum analogues of random walks, namely, quantum walks, have been developed for solving different problems. In 2003, Ambainis~\cite{r5} proposed a quantum walk algorithm for solving the element distinctness problem. His algorithm also solves the Markov chain search problem on the Johnson graph with cost $O(\mathsf{S} + \frac{1}{\sqrt{\varepsilon}}(\frac{1}{\sqrt{\delta}}\mathsf{U} + \mathsf{C}))$. In 2004, Szegedy~\cite{r45} gave a quantum walk algorithm for more generalized Markov chains with cost $O(\mathsf{S} + \frac{1}{\sqrt{\varepsilon\delta}}(\mathsf{U} + \mathsf{C}))$. We can view Szegedy's quantum walk as a quantum counterpart of a random walk, where one checks the state after each transition. Szegedy's quantum walk only detects the presence of a marked state, but cannot find one without extra costs. In 2006, Magniez et al.~\cite{r33} proposed a quantum walk search framework that unified the advantages of the quantum walks in~\cite{r5} and~\cite{r45}. In this quantum walk framework, we can perform the following operations:
\begin{itemize}
  \item \textbf{Setup:} with cost $\mathsf{S}$. preparing the initial state $\ket{\pi} = \frac{1}{\sqrt{|X|}}\sum_{x}\sqrt{\pi_x}\ket{x}$.
  \item \textbf{Update:}  with cost $\mathsf{U}$. applying the transformation $\ket{x}\ket{0} \mapsto \ket{x}\sum_{y\in X}\sqrt{p_{xy}}\ket{y}$.
  \item \textbf{Checking:} with cost $\mathsf{C}$, applying the transformation: $\ket{x} \mapsto \big\{\begin{smallmatrix} -\ket{x} \quad &\text{if $x\in M$} \\ \ket{x} \quad &\text{otherwise.}\end{smallmatrix}$
\end{itemize}
The main result of \cite{r33} is summarized as follows.
\begin{lemma}[\cite{r33}]
  \label{lem:framework}
  Let $P$ be an irreducible and ergodic Markov chain $P$ on $X$. Let $M \subseteq X$ be a subset of marked elements. Let $\varepsilon: = |M|/|X|$ and $\delta$ be the spectral gap of $P$. Then, there exists a quantum algorithm that with high probability, determines $M$ is empty or finds an $x \in M$ with cost $O(\mathsf{S} + \frac{1}{\sqrt{\varepsilon}}(\frac{1}{\sqrt{\delta}}\mathsf{U} + \mathsf{C}))$.
\end{lemma}

To solve the element distinctness problem, we define a Markov chain, following the work~\cite{r5,r11,r28}. The state space $X$ is all subsets of $[n]$ with size $r$. The Markov chain is based on the Johnson graph on $X$, where an edge is connecting $S$ and $S'$ if and only if $|S \cap S'| = r-1$. The transition probability on each edge is hence $\frac{1}{r(n-r)}$. A state $S$ is marked when there exist distinct $i, j \in S$ such the $i^{th}$ and the $j^{th}$ items are the same. The Markov chain has spectral gap $\delta \geq 1/r$ (see~\cite{r28}) and it is easy to verify that $\varepsilon \geq \binom{n-2}{r-2}/\binom{n}{r} = O(r^2/n^2)$. If we only consider the query complexity, the setup procedure costs $r$ queries, the update procedure costs one query, and the checking procedure does not cost any query. Choosing $r = n^{2/3}$ yields the optimal query complexity $O(n^{2/3})$.



\setcounter{section}{2}
\setcounter{theorem}{22}
\section{Quantum fine-grained complexity}

In this section, we give the formal definitions of the quantum fine-grained reduction and quantum strong exponential time hypothesis (QSETH). Moreover, we show that under QSETH, for $d=\polylog(n)$, the lower bounds for $\CP_{n,d}$ and $\OV_{n,d}$ are $n^{1-o(1)}$, which nearly matches the upper bounds given in \cref{cor:trivial_alg}.  

\subsection{Quantum fine-grained reduction and QSETH}

QSETH is defined based on the assumption that the best quantum algorithm for \textsf{CNF-SAT} is Grover search when the clause width $k$ is large enough. 
\begin{hypothesis}[QSETH]
For every $\epsilon > 0$, there exists a $k = k(\epsilon) \in \N$ such that no quantum algorithm can solve $k$-$\SAT$ (i.e., satisfiability on a CNF of width $k$) in $O(2^{(1/2 - \epsilon)n})$ time where $n$ is the number of variables. Moreover, this holds even when the number of clauses is at most $c(\epsilon) n$ where $c(\epsilon)$ denotes a constant that depends only on $\epsilon$.
\end{hypothesis}

Obviously, the Grover search can solve \textsf{CNF-SAT} in $\tilde{O}(2^{n/2})$. To the best of the our knowledge, there is no quantum algorithm that can do better than $O(2^{n/2})$ for any $k$. 


We recall that in the quantum query model, the input of a problem is given by a quantum oracle. Specifically, let $\mathsf{P}$ be a problem, and $X$ be an instance of $\mathsf{P}$ in the classical setting. Then, in the quantum query model, $X$ will be given by an oracle $\ora_X$. We will denote an algorithm or an oracle $\A$ with access to $\ora_X$ by $\A(\ora_X)$. 

We say $\A_{\epsilon}$ is an $\epsilon$-oracle for problem $\mathsf{P}$, if for every instance $\ora_X$, it holds that
\begin{align}
    \Pr[\A_{\epsilon}(\ora_X) = \mathsf{P}(X)] \geq 1-\epsilon,
\end{align}
and the running time is $O(1)$, where $\mathsf{P}(X)$ is the answer of $X$ for problem $\mathsf{P}$. 

\begin{definition}[Quantum oracles]\label{def:oisgood}
Let $X:=\{x_1,\dots,x_n\}$ be an instance of some problem and $\ora_X$ be the corresponding quantum oracle. To realize $\ora_X$, we do not need to write down the whole $X$; instead, we can just design a quantum circuit to realize the mapping 
\begin{align}
    \ket{i}\ket{0} \xrightarrow{\ora_X} \ket{i}\ket{x_i}.
\end{align}  
\end{definition}

\begin{definition}[Quantum fine-grained reduction]\label{def:qfgr}
Let $p(n)$ and $q(n)$ be nondecreasing functions of $n$. Let $\mathsf{P}$ and $\mathsf{Q}$ be two problems in the quantum query model and $\A_{\epsilon}$ be an $\epsilon$-oracle for $\mathsf{Q}$ with error probability $\epsilon\leq 1/3$. $\mathsf{P}$ is quantum ($p,q$)-reducible to $\mathsf{Q}$, denoted as $(\mathsf{P},p)\leq_{\QFG} (\mathsf{Q},q)$, if for every $\epsilon$, there exits a $\delta>0$, and algorithm $R$ with access to $\A_{\epsilon}$, a constant $d$, and an integer $k(n)$, such that for every $n\ge 1$, the algorithm $R$ takes any instance of $\mathsf{P}$ of size $n$ and satisfies the following: 
\begin{itemize}
    \item $R$ can solve $\mathsf{P}$ with success probability at least $2/3$ in time at most $d\cdot p(n)^{1-\delta}$.
    \item $R$ performs at most $k(n)$ quantum queries to $A_\epsilon$. Specifically, in the $j^{th}$ query, let $\mathbf{X}_j:= \{X_{1,j},X_{2,j},\dots\}$ be a set instances of $\mathsf{Q}$. Then, $R$ realizes the oracles $\{\ora_{X_{1,j}},\ora_{X_{2,j}},\dots\}$ in superposition and applies $A_{\eps}$ to solve the instances.
    \item The following inequality holds.
    \begin{align*}
        \sum_{j=1}^{k(n)}c(\mathbf{X}_j)\cdot q(n_j)^{1-\epsilon} \leq d\cdot p(n)^{1-\delta}, 
    \end{align*}
    where $c(\mathbf{X}_j)$ is the time required for $R$ to realize the oracles $\{\ora_{X_{1,j}},\ora_{X_{2,j}},\dots\}$ in superposition and $n_j := \max_{i}|X_{i,j}|$. 
\end{itemize}
\end{definition}

In \cref{def:qfgr}, the input of $A_{\eps}$ is given as a quantum oracle such that $A_{\eps}$ can be a quantum query algorithm with running time strictly less than the input size. Moreover, the quantum reduction $R$ can realize quantum oracles $\{\ora_{X_{1,j}},\ora_{X_{2,j}},\dots\}$ in superposition, and thus the time required is $\max_i c(X_{i,j})$ (where $c(X_{i,j})$ is the time required to realize $\ora_{X_{i,j}}$) instead of $\sum_{i} c(X_{i,j})$. This also allows $R$ to use fast quantum algorithms to process the information of $A_{\epsilon}'s$ output (e.g., amplitude amplification). 

\subsection{Lower bounds for $\CP$, $\OV$, and $\BCP$ in higher dimensions under QSETH}

Here, we give nearly linear lower bounds for $\OV$ and $\CP$ under QSETH by showing that there exist quantum fine-grained reductions from $\textsf{SAT}$ to these problems. 

\begin{theorem}\label{thm:qseth}
Assuming QSETH, for all $\epsilon>0$, there exists a $c$ such that $\OV_{n,c\log n}$ and $\CP_{n,(\log n)^c}$ cannot be solved by any quantum algorithm in time $O(n^{1-\epsilon})$.
\end{theorem}


We prove \cref{thm:qseth} by showing that there exist quantum fine-grained reductions from \textsf{CNF-SAT} to $\OV$, $\OV$ to $\BCP$, and $\BCP$ to $\CP$ with desired parameters.  We first give the reduction from \textsf{CNF-SAT} to $\OV$ as a warm-up. 

\begin{lemma}\label{lem:sat_q_ov}
\begin{align}
    (\mathsf{CNF\text{-}SAT}_{n},2^{n/2})\leq_{\QFG} (\OV_{n_1,d_1},n_1),
\end{align}
where $n_1=2^{n/2}$ and $d_1=\Theta(n)$.
\end{lemma}
\begin{proof}
Let $\phi$ be a CNF formula with $n$ variables and $m=\Theta(n)$ clauses. Let $\mathcal{A}$ be an algorithm for $\OV$. We first recall the classical reduction. Let $\phi := \phi_1\wedge \cdots \wedge\phi_m$. 
We divide the $n$ variables into two sets $A$ and $B$ with $|A| = |B| = \frac{n}{2}$. Let $A:=\{x_1,\dots,x_{n/2}\}$ and $B:= \{x_{n/2+1},\dots,x_{n}\}$. We let $S_A:=\{a_1,\dots,a_{2^{n/2}}\}$ be all assignments to $A$ and $S_B:=\{b_1,\dots,b_{2^{n/2}}\}$ be all assignments to $B$. We describe two mappings $f_A: S_A\rightarrow \{0,1\}^m$ and $f_B:S_B\rightarrow \{0,1\}^m$ as follows: 
\begin{align}
    &f_A(a_i) = [\phi_1(a_i),\dots,\phi_m(a_i)]^T, \text{ and}\\
    &f_B(b_i) = [\phi_1(b_i),\dots,\phi_m(b_i)]^T,
\end{align}
where $\phi_j(a_i) = 0$ if $a_i$ is a satisfied assignment for $\phi_j$, and  $\phi_j(a_i) = 1$ otherwise; we define $\phi_i(b_i)$ in the same way. Let $F_A:= \{f_A(a_i):\; i\in[2^{n/2}]\}$ and $F_B:= \{f_B(b_i):\; i\in[2^{n/2}]\}$. Then, it is obvious that if there exist $v\in F_A$ and $u\in F_B$ such that $\langle v, u \rangle=0$, then $\phi$ is satisfiable. However, at first glance, this reduction with $O(2^{n/2})$ running time is not fine-grained since we require the cost of the reduction to be at most $2^{n(1-\delta)/2}$ for some $\delta>0$ by \cref{def:qfgr}, but writing down elements in $F_A$ and $F_B$ already takes $\Omega(2^{n/2})$. %

Nevertheless, as in \cref{def:oisgood}, a quantum fine-grained reduction only needs to realize the functions $f_A$ and $f_B$, which takes $O(mkn)$ time where $k$ is the width of clauses. This is much less than $O(2^{n(1-\delta)/2})$. More specifically,  $f_A$ and $f_B$ are oracles for $F_A$ and $F_B$, and for any quantum query to elements in $F_A$ or $F_B$, the reduction can implement oracles $f_A$ and $f_B$:   
\begin{align}
    \ket{e,x}\ket{0}\xrightarrow{f_e} \ket{e,x}\ket{f_e(x)},
\end{align}
where $e\in \{A,B\}$, and the time $c(f_e)$ for the reduction to implement $f_e$ for one quantum query is at most $O(kmn)$. Finally, this reduction only uses one oracle ($F_A,F_B$). If there is an algorithm for $\OV$ which succeeds with probability $2/3$, we can boost the success probability of the reduction by repetition.  Therefore, ($\textsf{CNF-SAT},2^{n/2}$) is quantum reducible to ($\textsf{OV}_{n_1,d_1},n_1$). 
\end{proof}

Then, to prove $(\textsf{CNF-SAT}, 2^{n/2})\leq_{\QFG} (\textsf{CP}_{n_3,d_3}, n_3)$, we show that $(\textsf{BCP}_{n_2,d_2},n_2)\leq_{\QFG} (\CP_{n_3,d_3},n_3)$ and $(\textsf{OV}_{n_1,d_1}, n_1) \leq_{\QFG} (\textsf{BCP}_{n_2,d_2},n_2)$, where $n_2, n_3, d_2,d_3$ are some functions of $n$ specified in the following lemmas.

\begin{lemma}\label{lem:qbcp_to_cp}
For $d=\Theta(\log n)$,
\begin{align}
    (\BCP_{n,d},n)\leq_{\QFG} (\CP_{n',d'},n'),
\end{align}
where $n' = n^{O(1)}$ and $d'=(\log n)^{c}$ for some constant $c$ and all points have $\{0,1\}$ entries with the Hamming metric.
\end{lemma}

\begin{remark}
The points have coordinate entries in $\{0,1\}$, and the Hamming metric is equivalent to distance in $\ell_2$-metric (up to power of 2) in this case. Therefore, in the proof of Lemma~\ref{lem:qbcp_to_cp}, we can consider the Hamming distance between points instead of $\ell_2$ distance without loss of generality. 
\end{remark}

We first introduce the classical reductions in~\cite{r29} and some results we will use to prove Lemma~\ref{lem:qbcp_to_cp}. 

\paragraph{Classical reduction} 
We can consider an instance of $\BCP$ with two sets of points $A$ and $B$ as a weighted complete bipartite graph $K_{n,n}$, where the vertices are the points in these two sets and edges' weights are equal to the distances between the corresponding points. Then, solving $\BCP$ is equivalent to find an edge with the minimum weight in this graph.  
However, we cannot directly apply the algorithm for $\CP$ on this graph since there could be two points in the same set (no edge connecting them) that have a smaller distance than any pairs of points in two sets (connected by an edge). To overcome this difficulty, we can ``stretch'' the points to make the points in the same set far from each other, which is characterized by the contact dimension of a graph:
\begin{definition}[Contact Dimension]\label{def:contact_dim}
For any graph $G=(V,E)$, a mapping $\tau:V\rightarrow \R^d$ is said to realize $G$ if for some $\beta>0$, the following holds for every distinct vertices $u,v$:
\begin{align}
    \|\tau(u)-\tau(v)\|_2  =&~\beta ~ \text{if~} \{u,v\}\in E,\\
    \|\tau(u)-\tau(v)\|_2  > &~\beta ~ \text{otherwise.}\notag
\end{align}
The contact dimension of $G$, denoted by $\mathrm{cd}(G)$, is the minimum $d\in \N$ such that there exists $\tau:V\rightarrow \R^d$ realizing $G$.
\end{definition}
That is, with the help of $\tau$, we can restrict the optimal solution of $\CP$ to be the points connected by an edge in $G$. But we cannot realize the whole complete bipartite graph since $\mathrm{cd}(K_{n,n})=\Theta(n)$, which makes the dimension of the $\CP$ instance too large. \cite{r29} showed that we can realize a subgraph of $K_{n,n}$ and apply permutations to its vertices such that the union of these subgraphs cover $K_{n,n}$. In this way, $\BCP$ can be computed by solving $\CP$ on each subgraph and outputting the best solution. More specifically, the reduction in \cite{r29} relies on the following theorem:

\begin{theorem}[Theorem 4.2 in \cite{r29}] \label{thm:dense-cd}
For every $0<\delta < 1$, there exists a log-dense sequence $(n_i)_{i\in \N}$ such that, for every $i\in \N$, there is a bipartite graph $G_i=(A_i \dot \cup B_i, E_i)$ where $|A_i|=|B_i|=n_i$ and $|E_i|\ge \Omega(n_i^{2-\delta})$, such that $\mathrm{cd}(G_i)=(\log n_i)^{O(1/\delta)}$. Moreover, for all $i\in \N$, a realization $\tau:A_i\dot \cup B_i\rightarrow \{0,1\}^{(\log n_i)^{O(1/\delta)}}$ of $G_i$ can be constructed in $n_i^{2+o(1)}$ time.
\end{theorem}
\noindent The log-dense sequence is defined as follows:
\begin{definition}
A sequence $(n_i)_{i\in \N}$ of increasing positive integers is log-dense if there exists a constant $c\ge 1$ such that $\log n_{i+1}\le c\cdot \log n_i$ for all $i\in \N$.
\end{definition}
They also showed that, the permutations for covering the complete bipartite graph can be efficiently found, as shown in the following lemma.
\begin{lemma}[Lemma 3.11 in \cite{r29}] \label{lem:cover}
For any bipartite graph $G(A\dot \cup B, E_G)$ where $|A|=|B|=n$ and $E_G\ne \emptyset$, there exist side-preserving permutations $\pi_1,\dots, \pi_k:A\cup B\rightarrow A\cup B$ where $k\leq \frac{2n^2\ln n}{|E_G|}+1$ such that
\begin{align}
    \underset{i\in [k]}{\bigcup} E_{G_{\pi_i}} = E_{K_{n,n}}.
\end{align}
Moreover, such permutations can be found in $O(n^6\log n)$ time.
\end{lemma}

Now, we are ready to state the quantum fine-grained reduction by ``quantizing'' the classical reduction.
\begin{proof}[Proof of \cref{lem:qbcp_to_cp}]
Let $A,B$ be the two sets of input points of $\textsf{BCP}$. Suppose for \textsf{BCP}, there is an input oracle $\ora_{\BCP}$ which, given an index, returns the corresponding point:
\begin{align}
    \ket{b}\ket{i}\ket{0} \xrightarrow{\ora_{\BCP}}
    \begin{cases}
        \ket{b}\ket{i}\ket{x_i} &\text{if } b= 0,\\
        \ket{b}\ket{i}\ket{y_i} & \text{if } b= 1,
    \end{cases}
\end{align}
where $x_i$ is the $i$-th point in the set $A$ and $y_i$ is the $i$-th point in the set $B$. The sizes of $A$ and $B$ are both equal to $n$ and each point is in $\{0,1\}^{d_1}$, where $d_1=\Theta(\log n)$ is the dimension of $\BCP$. 
 
For $\CP$, suppose there is a quantum algorithm $\mathcal{A}$ such that for $m$ points in $\{0,1\}^{d_2}$ given by an oracle $\mathcal{M}_{CP}$, $\mathcal{A}^{\mathcal{M}_{CP}}$ returns the closest pair of these $n$ points with probability at least $2/3$.

Then we need to transform $\ora_{\BCP}$ to some oracles $\mathcal{M}_i$ for \CP, such that by running $\mathcal{A}$ with $\mathcal{M}_i$ as input oracles, we can get the bichromatic closest pair between $A$ and $B$. The reduction has four steps:

\paragraph*{1. Pre-processing.} We first follow the classical reduction to pre-process the input points of \textsf{BCP}. For some integer $n'\leq n^{0.1}$, we can partition $A$ and $B$ into $n'$-size subsets:
\begin{align}
    A= &~ A_1~\dot\cup ~\cdots ~\dot \cup~ A_{r},\\
    B= &~ B_1 ~\dot\cup ~\cdots ~\dot \cup~ B_{r},\notag
\end{align}
where $r=\lfloor n/n'\rfloor$. Here, we assume that $n$ is divisible by $n'$. It follows that
\begin{align}\label{eq:bcp_bcp}
    \textsf{BCP}(A,B)=\min_{i,j\in [r]} \textsf{BCP}(A_i,B_j).
\end{align}

Then, we use the algorithm in \cite{r29} to construct $k$ mappings $f_1,\dots, f_k:[2n']\rightarrow \{0,1\}^{d'}$ such that 
\begin{align}\label{eq:bcp_cp}
    \textsf{BCP}(A_i,B_j)=\min_{t\in [k]} \textsf{CP}(f_t(A_i)\cup f_t(B_j)) ~~~ \forall i,j\in [\lfloor n/n'\rfloor].
\end{align}

More specifically, we pick $n'$ to be the largest number in a log-dense sequence that is smaller than $n^{0.1}$. Then, we apply~\cref{thm:dense-cd} to classically construct a bipartite graph $G(A\cup B,E)$ with $n'$ vertices in each side and a realization $\tau$.
By choosing $\delta = \epsilon/2$ in \cref{thm:dense-cd}, the graph $G$ has $|E|=\Omega(n'^{2-\epsilon/2})$ edges. And we can get $2n'$ 0/1-strings of length $(\log n')^{O(2/\epsilon)}$:
\begin{align}
    \tau^A_i=\tau(u_i)~~~\forall u_i\in A, \quad\text{ and }\quad
    \tau^B_i=\tau(v_i)~~~\forall v_i\in B.
\end{align}

In order to cover the complete bipartite graph, we run the classical algorithm (\cref{lem:cover}) to find $k$ permutations $\pi_1,\dots, \pi_k:[n']\rightarrow [n']$, where $k$ is a parameter to be chosen later.

Then, we can define the mappings as follows:  
\begin{align}\label{eq:fcn_cp}
    f_t(u)=
    \begin{cases}
        x_v\circ \left(\tau^A_{\pi_t(w)}\right)^{d+1} & \text{if }1\leq u\leq n'\\
        y_v\circ \left(\tau^B_{\pi_t(w)}\right)^{d+1} & \text{if }n'<u\leq 2n'
    \end{cases} ~~~\forall t\in [k], u\in [2n'],
\end{align}
where $\circ$ means string concatenation and $(s)^{d+1}$ denotes $d+1$ copies of the string $s$. For a point $p\in A_i\cup B_j$, $u\in [2n']$ is the index in this union-set,  $v\in [n]$ is the index in the ground set $A$ or $B$, and $w\in [n']$ is the index in the subset $A_i$ or $B_j$. Further, if $1\leq u\leq n'$, then $w:=u$; otherwise, $w:=u-n'$.

\paragraph{2. Oracle construction.}
For $i,j\in [r],t\in [k]$, we then construct the input oracle $\mathcal{M}_{i,j,t}$ for the problem $\textsf{CP}(f_t(A_i)\cup f_t(B_j))$. For a query index $u\in [2n']$, 
\begin{align}\label{eq:oracle_cp}
    M_{i,j,t}\ket{u}\ket{0}=\ket{u}\ket{f_t(u)}.
\end{align}

With the help of the input oracle $\ora_{\BCP}$, we can implement $\mathcal{M}_{i,j,t}$ in the following way: 
\begin{enumerate}
    \item Prepare an ancilla qubit $\ket{b}$ such that $b=1$ if $u>n'$. 
    \item Transform $\ket{u}$ to $\ket{v}$, the index of the point in $A$ or $B$, based on the value of $b$. Note that the index is unique. Hence, this transformation is unitary and can be easily achieved by a small quantum circuit.
    \item Query $\ora_{\BCP}$ with input $\ket{b}\ket{v}$. Assume $b=0$. Then, 
    \begin{align}
        \ket{b}\ket{v}\ket{0}\xmapsto{\ora_{\BCP}}\ket{b}\ket{v}\ket{x_v}.
    \end{align}
    \item Similar to the second step, the index $w$ of the point in $A_i$ and $B_j$ can be computed from $v$ by a unitary transformation:
    \begin{align}
        \ket{b}\ket{v}\ket{x_v}\mapsto \ket{b}\ket{w}\ket{x_v}
    \end{align}
    \item Since each $w$ corresponds to a unique string $\tau^A_{\pi_t(w)}$, we can attach $d+1$ copies of this string to the remaining quantum registers:
    \begin{align}
        \ket{b}\ket{w}\ket{x_v}\mapsto \ket{b}\ket{w}\ket{x_v}\ket{\left( \tau^A_{\pi_t(w)} \right)^{d+1}}.
    \end{align}
    \item By recovering $u$ from $w$, we get the final state:
    \begin{align}
        \ket{u}\ket{f_t(u)}=\ket{u}\ket{x_v,\left( \tau^A_{\pi_t(w)} \right)^{d+1}}.
    \end{align}
\end{enumerate}

\paragraph*{3. Query process}
By \cref{eq:bcp_bcp,eq:bcp_cp}, we have
\begin{align}\label{eq:bcp_all}
    \textsf{BCP}(A,B)=\min_{i,j\in [r],t\in [k]}\textsf{CP}(f_t(A_i)\cup f_t(B_j)).
\end{align}
Hence, we can use quantum minimum-finding algorithm in \cref{cor:trivial_alg} over the sub-problems to find the minimum solution. For each sub-problem, we can run the algorithm for \textsf{CP} with $\mathcal{M}_{i,j,t}$ as the input oracle.

\paragraph*{4. Post-processing.} In case that $n$ is not divisible by $n'$, let the remaining points in $A$ and $B$ be $A_{res}$, $B_{res}$, respectively. Then, we can use Grover search to find the closest pair between $A_{res}$ and $B$, and between $B_{res}$ and $A$. Then, compare the answer to the previously computed result and pick the smaller one.

\paragraph*{Correctness.} In this reduction, we do not change the constructions of the mappings $\{f_i\}_{i\in [k]}$. By~\cite{r29}, \cref{eq:bcp_all} is correct in the classical setting. Hence, it also holds in the quantum setting, and we can use Grover search to find the minimum solution. However, since the algorithm $\mathcal{A}$ for $\CP$  has success probability $2/3$, for each tuple $(i,j,t)\in [r]\times [r]\times [k]$, we need to run $\mathcal{A}^{\mathcal{M}_{i,j,t}}$ $O(\log n)$ times to boost the success probability to at least $1-\frac{1}{n}$. Then, by the union bound, the probability that all queries in the Grover search are correct is at least $99/100$. Hence, by \cref{thm:mim}, the overall success probability is at least $2/3$.  


\paragraph*{Running Time of the Reduction.} 

The running time of the pre-processing step consists of two parts: (1) constructing the graph $G$ and its realization $\tau$; (2) finding $k$ permutations. For the first part, by~\cref{thm:dense-cd}, it can be done in $n'^{2+o(1)}$ time. For the second part, we pick $k=O(\frac{2n'^2\log n'}{n'^{2-\epsilon/2}})=O(n'^{\epsilon/2}\log n')$, and by~\cref{lem:cover}, it can be done in $O(n'^6 \log n')$ time. Hence, the total running time of pre-processing step is $n'^{2+o(1)}+O(n'^6\log n')=\wt{O}(n^{0.6})$.

The oracle construction can be done ``on-the-fly''. More specifically, given the strings $\{\tau^A_i,\tau^B_i\}_{i\in [n']}$, and permutations $\{\pi_i\}_{i\in [k]}$, for each query index $u$, we can simulate the oracle $\mathcal{M}_{i,j,t}$ defined in \cref{eq:oracle_cp} in $c(\mathcal{M}_{i,j,t})=O(d_2)=(\log n')^{\Omega(1)}=\wt{O}(1)$ time.

In the query process, for each $\CP$ instance indexed by $(i,j,t)$, suppose $\mathcal{A}^{\mathcal{M}_{i,j,t}}$ gets the answer in time $q(n')=n'$. Moreover, for each time $\mathcal{A}$ querying the input oracle $\mathcal{M}_{i,j,t}$, we need to spend $c(\mathcal{M}_{i,j,t})$ time to simulate the oracle. And we also have $O(\log n)$ runs for each instance. Hence, the total running time for each $\CP$ is at most  
\begin{align}
    n'^{1-\epsilon}\cdot \wt{O}(1)\cdot O(\log n)=\wt{O}(n'^{1-\epsilon}).
\end{align}
Then, we use Grover's search algorithm over $r^2\cdot k$ instances, which can be done by querying $\wt{O}(\sqrt{r^2\cdot k})$ instances by \cref{thm:mim}. Therefore, for any $\epsilon > 0$, we have

\begin{align}\label{eq:bcp_qfgr}
    \wt{O}(\sqrt{r^2 k})~\cdot &~ q(n')^{1-\epsilon}\cdot c(\mathcal{M}_{i,j,t})\cdot O(\log n) =  ~ 
    \wt{O}(\sqrt{(n/n')^2 k}\cdot(n')^{1-\eps}) \\
    \leq &~ \wt{O}(n\cdot (n')^{-\eps})
    \leq \wt{O}(n\cdot n^{-\eps/2})
    \leq n^{1-\delta},
\end{align}
where the first inequality follows from $k=O(n'^{\epsilon / 2}\log n')$ as shown in \cite{r29} and the last inequality follows by setting $\delta = \epsilon / 10$.

For the post-processing step, the sizes of $A_{res}$ and $B_{res}$ are at most $n'$. The running time is 
\begin{align}
    O(\sqrt{n\cdot n'}\cdot \log n) \leq \wt{O}(n^{0.55}).
\end{align}

Therefore, for any $\epsilon>0$, there exists a $\delta>0$ such that the \cref{eq:bcp_qfgr} holds and the total reduction time is $O(n^{1-\delta})$. By \cref{def:qfgr}, $\BCP_{n,d_1}$ can be quantum fine-grained reduced to $\CP_{n,d_2}$. This completes the proof of this lemma. 
\end{proof}

Finally, we show that $(\textsf{OV}_{n,d},n)\leq_{\QFG} (\textsf{BCP}_{n,d'},n)$ by quantizing the reduction in~\cite{r29} following the same idea. 
\begin{lemma}\label{lem:ov_q_bcp}
For $d=\Theta(\log n)$, 
\begin{align} 
    (\OV_{n,d},n)\leq_{\QFG} (\BCP_{n,d'},n),
\end{align}
where $d'=\Theta(\log n)$.
\end{lemma}
\begin{proof}
For an $\OV$ instance with sets of vectors $A$ and $B$, let $\ora_{\OV}$ be the input oracle such that:
\begin{align}
    \ora_{\OV}\ket{i}\ket{0}=\begin{cases}
        \ket{i}\ket{a_i} & \text{if }i\in A,\\
        \ket{i}\ket{b_i} & \text{if }i\in B.
    \end{cases}
\end{align}
where $a_i,b_i\in \{0,1\}^d$.

Then, similar to the classical reduction, we can construct mappings $f_A, f_B:\{0,1\}^d\rightarrow \{0,1\}^{5d}$ such that
\begin{align}
    f_A(a_i)_{5j-4:5j}=\begin{cases}
        11000 & \text{if }a_i(j)=0\\
        00110 & \text{if }a_i(j)=1
    \end{cases}~~~\forall j\in [d],
\end{align}
and
\begin{align}
    f_B(b_i)_{5j-4:5j}=\begin{cases}
        10100 & \text{if }b_i(j)=0,\\
        01001 & \text{if }b_i(j)=1.
    \end{cases} ~~~\forall j\in [d].
\end{align}

By the classical reduction, we have
\begin{align}
    \textsf{OV}(A,B)=1 \text{ if and only if }\textsf{BCP}(f_A(A),f_B(B))=2d
\end{align}
under Hamming distance. 

Also, note that we can simulate the input oracle $\ora_{\BCP}$ by first querying the oracle $\ora_{\OV}$ to get the vector, then applying the corresponding mapping $f_A$ or $f_B$, which can be done in $c(\ora_{\BCP})=O(d)$ time. Let the running time of the algorithm for $\textsf{BCP}$ be $q(n)=n$. Then for any $\epsilon>0$,
\begin{align}
    q(n)^{1-\epsilon}\cdot c(\ora_{\BCP})=n^{1-\epsilon}\cdot \Theta(\log n)\leq n^{1-\delta}
\end{align}
for some small $\delta >0$. Hence, by \cref{def:qfgr}, $(\textsf{OV}_{n,d},n)\leq_{\QFG} (\textsf{BCP}_{n,d'},n)$.
\end{proof}
 
\begin{proof}[Proof of \cref{thm:qseth}]

We can prove the theorem by contradiction following \cref{lem:sat_q_ov}, \cref{lem:ov_q_bcp}, and \cref{lem:qbcp_to_cp}. Specifically, suppose that there exists an $\epsilon>0$, for all $d=\Theta(\log n)$, there exists a quantum algorithm which can solve $\OV$ in time $O(n^{1-\epsilon})$. Then, we can obtain a quantum algorithm for \textsf{CNF-SAT}, which runs in time $O(2^{n/2(1-\epsilon)})$ by \cref{lem:sat_q_ov}. This contradicts QSETH. The proof for $\CP$ is the same. 
\end{proof}

\subsection{Quantum lower bound for $\BCP$ in nearly-constant dimensions under QSETH}
A byproduct of the previous subsection is a quantum lower bound for $\BCP$ in higher dimensions (i.e., $d = \polylog(n)$) under QSETH (\cref{lem:ov_q_bcp}). In this subsection, we show that this quantum lower bound for $\BCP$ even holds for nearly-constant dimensions (i.e., $d = c^{\log^*(n)}$). The main result of this subsection is the following theorem.
\begin{theorem}\label{thm:qseth_2_bcp}
Assuming QSETH, there is a constant $c$ such that $\BCP$ in $c^{\log^* (n)}$ dimensions requires $n^{1-o(1)}$ time for any quantum algorithm.
\end{theorem}

We will ``quantize'' the results by Chen~\cite{r17} to prove this theorem. More specifically, we first show a quantum fine-grained self-reduction of $\OV$ from $\log n$ dimensions with binary entries to $2^{\log ^* (n)}$ dimensions with integer entries ($\ZOV$). Then, we give a quantum fine-grained reduction from $\ZOV$ to $\BCP$ in nearly-constant dimensions.

\begin{definition}[Integral Orthogonal Vector, $\ZOV$]
Given two sets $A,B$ of $n$ vectors in $\mathbb \Z^d$, find a pair of vectors $a\in A$ and $b\in B$ such that
$\ip{a}{b}=0$, where the inner product is taken in $\mathbb{Z}$.
\end{definition}

We use $\ZOV_{n,d}$ to denote $\ZOV$  with $n$ vectors of $d$ dimension in each set. We then recap a theorem in \cite{r17}:

\begin{theorem}[{\cite[Theorem 4.1]{r17}}]\label{thm:zov_map}
Let $b,\ell$ be two sufficiently large integers. There is a classical reduction $\psi_{b,\ell}:\{0,1\}^{b\cdot\ell}\rightarrow \Z^\ell$ and a set $V_{b,\ell}\subseteq \Z$, such that for every $x,y\in \{0,1\}^{b\cdot \ell}$,
\begin{align}\label{eq:zov_iff}
    \ip{x}{y}=0~\Leftrightarrow~\ip{\psi_{b,\ell}(x)}{\psi_{b,\ell}(y)}\in V_{b,\ell}
\end{align}
and
\begin{align}
    0\leq \psi_{b,\ell}(x)_i< \ell^{6^{\log^* (b)}\cdot b}
\end{align}
for all possible $x$ and $i\in [\ell]$. Moreover, the computation of $\psi_{b,\ell}(x)$ takes $\poly(b\cdot \ell)$ time, and the set $V_{b,\ell}$ can be constructed in $O\left( \ell^{O(6^{\log^* (b)}\cdot b)}\cdot \poly(b\cdot \ell) \right)$ time.
\end{theorem}

Note that the size of $V_{b,\ell}$ is at most $\ell^{2\cdot 6^{\log^* (b)}\cdot b+1}$. 
The following lemma gives a quantum fine-grained reduction from $\OV$ to $\ZOV$:
\begin{lemma}\label{lem:ov_2_zov}
For $d=\Theta(\log n)$,
\begin{align}
    (\OV_{n,d}, n)\leq_{\QFG} (\ZOV_{n,d'}, n).
\end{align}
where $d'=2^{O(\log^* n_2)}$.
\end{lemma}
\begin{proof}
Consider an $\OV_{n,d}$ with $d=c\cdot \log n$, where $c$ is an arbitrary constant. We choose $\ell:=7^{\log^* n}$ and $b:=d/\ell$. Then, we can apply~\cref{thm:zov_map} to get the mapping function $\psi_{b,\ell}$ and the set $V_{b,\ell}$. For each $v\in V_{b,\ell}$, we'll construct an instance of $\ZOV_{n,\ell+1}$ as follows:
\begin{enumerate}
    \item Let $\ket{i}$ be the input query index of $\ZOV_{n,\ell+1}$.  
    \item Query $\OV_{n,d}$'s input oracle $\ora_{\OV}$ and get the vector $\ket{i,x}$.
    \item Compute the mapping $\psi_{b,\ell}$ and get $\ket{i,x}\ket{\psi_{b,\ell}(x)}$.
    \item If $x\in A$, then attach 1 to the end of the register: $\ket{i,x}\ket{\psi_{b,\ell}(x),1}$. If $x\in B$, then attach $-v$ to the end: $\ket{i,x}\ket{\psi_{b,\ell}(x),-v}$.
    \item Use $\ora_{\OV}$ to erase $x$ and return the final input state $\ket{i}\ket{\psi_{b,\ell}(x),1}$ or $\ket{i}\ket{\psi_{b,\ell}(x),-v}$.
\end{enumerate}
For each instance, we can use the quantum oracle for $\ZOV_{n,\ell+1}$ to check the orthogonality. $\OV_{n,d}$ is YES if and only if there exists a YES-instance of $\ZOV_{n,\ell+1}$.

\paragraph*{Correctness.} The correctness follows from~\cref{eq:zov_iff}: 
\begin{align}
\ip{x}{y}=0 \Leftrightarrow \ip{\psi_{b,\ell}(x)}{\psi_{b,\ell}(y)}=v\in V_{b,\ell} \Leftrightarrow \ip{[\psi_{b , \ell}(x), ~ 1]}{[\psi_{b,\ell}(y), ~ -v]}=0. 
\end{align}

\paragraph*{Reduction time.} Note that for $\ell=7^{\log^* n}$ and $b=d/\ell$, we have:
\begin{align}
    \log\left( \ell^{O(6^{\log^* (d)}\cdot b)} \right) = &~ \log \ell \cdot O\left( 6^{\log^* (d)}\cdot (d/\ell) \right)\\
    =&~ O\left( \log^* (n)\cdot 6^{\log^* n} \cdot c\log n / 7^{\log^* n} \right)\\
    =&~ o(\log n).
\end{align}
This implies that $|V_{b,\ell}|\leq \ell^{2\cdot 6^{\log^* (b)}\cdot b+1} \leq n^{o(1)}$.
Hence, the number of $\ZOV_{n,\ell+1}$ instances is $n^{o(1)}$ and the running time for compute $V_{b,\ell}$ is $n^{o(1)}$. And for each input query, the oracle for $\ZOV_{n,\ell+1}$ can be simulated in $c(\ora_{\ZOV})=\poly(d)=\poly(\log n)$ time. We can show that for every $\epsilon>0$, if $\ZOV_{n,\ell+1}$ can be decided in $n^{1-\epsilon}$ time, then
\begin{align}
    \sum_{v\in V_{b,\ell}} n^{1-\epsilon}\cdot c(\ora_{\ZOV})=n^{o(1)}\cdot n^{1-\epsilon}\cdot \poly(\log n) \leq n^{1-\delta}
\end{align}
for some $\delta > 0$, which satisfies the definition of quantum fine-grained reduction (\cref{def:qfgr}).

Therefore, $\OV_{n,O(\log n)}$ is quantum fine-grained reducible to $\ZOV_{n,2^{O(\log^* (n))}}$.
\end{proof}

Then, we give a quantum fine-grained reduction from $\ZOV$ to $\BCP$:
\begin{lemma}\label{lem:zov_2_bcp}
For $d=2^{O(\log^* n)}$,
\begin{align}
    (\ZOV_{n,d}, n)\leq_{\QFG} (\BCP_{n,d'}).
\end{align}
where $d'=d^2+2$.
\end{lemma}
\begin{proof}
We remark here that this proof closely follows that for Theorem 4.3 in \cite{r17}. We nonetheless give it here as some details are different.

For an $\ZOV_{n,d}$ instance with $(k\cdot \log n)$-bit entries, we construct a $\BCP$ instance as follows:
\begin{enumerate}
    \item For $x\in A$, construct a vector $x'\in \Z^{d^2}$ such that $x'_{i,j}=x_i\cdot x_j$. Here, we index a $d^2$-dimensional vector by $[d]\times [d]$. Similarly, for $y\in B$, construct a vector $y'\in \Z^{d^2}$ such that $y'_{i,j}=-y_i \cdot y_j$.
    \item Choose $W:=(d^2+1)\cdot n^{4k}$. For each $x'$, construct a vector $x''\in \R^{d^2+2}$ such that 
    \begin{align}
        x''=\left[x',~ \sqrt{W-\|x'\|_2^2},~ 0\right].
    \end{align}  For each $y'$, construct a vector $y''\in \R^{d^2+2}$ such that 
    \begin{align}
        y''=\left[y',~ 0, ~ \sqrt{W-\|y'\|_2^2}\right].
    \end{align}
\end{enumerate}
Then, we claim that the $\ZOV$ instance is YES if and only if the $\BCP$ instance has the minimum distance $\leq \sqrt{2W}$.

\paragraph*{Correctness.} First note that $\|x'\|_2^2\leq d^2\cdot (2^{k\log n})^4=d^2\cdot n^{4k}$. Hence, $W-\|x'\|_2^2 > 0$ and $W-\|y'\|_2^2>0$. For any $x''$ and $y''$ in the new constructed instance of $\BCP$, we have
\begin{align}
    \|x''-y''\|_2^2=&~ \|x''\|_2^2+\|y''\|_2^2-2\cdot \ip{x''}{y''}\\
    =&~ 2\cdot W-2\cdot \ip{x'}{y'}\\
    =&~ 2\cdot W - 2\cdot \sum_{(i,j)\in [d]\times [d]}x_i\cdot x_j\cdot (-y_j\cdot y_j)\\
    =&~ 2\cdot W + 2\cdot (\ip{x}{y})^2.
\end{align}
Hence,
\begin{align}
    \ip{x}{y}=0~\Leftrightarrow~\|x''-y''\|_2^2= 2W.
\end{align}

\paragraph*{Reduction time.} We can see from the above description that the input mapping function is simple and can be computed by a small quantum circuit in $O(d^2)=O(2^{O(\log^* (n))})$ time. Hence, we have $c(\ora_{\BCP})=O(2^{O(\log^* (n))})$. Also, by~\cref{def:qfgr}, it's easy to check that this is indeed a quantum fine-grained reduction from $\ZOV$ to $\BCP$.
\end{proof}

Now \cref{thm:qseth_2_bcp} follows immediately from \cref{lem:ov_2_zov} and \cref{lem:zov_2_bcp}:
\begin{proof}[Proof of \cref{thm:qseth_2_bcp}]
Let $\eps > 0$ be some constant. Suppose we can solve $\BCP_{n,c^{\log^* (n)}}$ in $n^{1-\epsilon}$ time for all constant $c>0$. Then, by \cref{lem:ov_2_zov} and \cref{lem:zov_2_bcp}, we can also solve $\OV_{n,c'\log n}$ in $n^{1-\delta}$ time for some $\delta > 0$ and any $c'>0$. However, this contradicts QSETH by~\cref{thm:qseth}. Therefore, assuming QSETH, there exists a constant $c$ such that $\BCP_{n,c^{\log^*( n)}}$ requires $n^{1-o(1)}$ time.
\end{proof}

\begin{thebibliography}{52}
\bibitem{r1} Scott Aaronson and Yaoyun Shi. Quantum Lower Bounds for the Collision and the Element Distinctness Problems. J. ACM, 51(4):595–605, July 2004.
\bibitem{r2} Amir Abboud, Ryan Williams, and Huacheng Yu. More Applications of the Polynomial Method to Algorithm Design. In Proceedings of the Twenty-sixth Annual ACM-SIAM Symposium on Discrete Algorithms, SODA ’15, pages 218–230, 2015.
\bibitem{r3} Pankaj K. Agarwal, Herbert Edelsbrunner, Otfried Schwarzkopf, and Emo Welzl. Euclidean Minimum Spanning Trees and Bichromatic Closest Pairs. Discrete \& Computational Geometry, 6(3):407–422, September 1991.
\bibitem{r4} A. Ambainis. Quantum Search Algorithms. SIGACT News, 35(2):22–35, June 2004. doi: 10.1145/992287.992296.
\bibitem{r5} Andris Ambainis. Quantum Walk Algorithm for Element Distinctness. SIAM Journal on Computing, 37(1):210–239, 2007.
\bibitem{r6} Sanjeev Arora and Boaz Barak. Computational Complexity: A Modern Approach. Cambridge University Press, 2009.
\bibitem{r7} Arturs Backurs, Piotr Indyk, and Ludwig Schmidt. On the Fine-Grained Complexity of Empirical Risk Minimization: Kernel Methods and Neural Networks. In Advances in Neural Information Processing Systems, pages 4308–4318, 2017.
\bibitem{r8} Marshall Ball, Alon Rosen, Manuel Sabin, and Prashant Nalini Vasudevan. Proofs of Useful Work. IACR Cryptology ePrint Archive, 2017:203, 2017.
\bibitem{r9} C. Bennett, E. Bernstein, G. Brassard, and U. Vazirani. Strengths and Weaknesses of Quantum Computing. SIAM Journal on Computing, 26(5):1510–1523, 1997.
\bibitem{r10} Jon Louis Bentley and Michael Ian Shamos. Divide-and-Conquer in Multidimensional Space. In Proceedings of the eighth annual ACM symposium on Theory of computing, pages 220–230. ACM, 1976.
\bibitem{r11} Daniel J. Bernstein, Stacey Jeffery, Tanja Lange, and Alexander Meurer. Quantum Algorithms for the Subset-Sum Problem. In Post-Quantum Cryptography, pages 16–33. Springer Berlin Heidelberg, 2013. doi:10.1007/978-3-642-38616-9\_2.
\bibitem{r12} Dominic W. Berry, Andrew M. Childs, Richard Cleve, Robin Kothari, and Rolando D. Somma. Exponential Improvement In Precision for Simulating Sparse Hamiltonians. Forum of Mathematics, Sigma, 5, 2017. doi:10.1017/fms.2017.2.
\bibitem{r13} Sergei N Bespamyatnikh. An Optimal Algorithm for Closest-Pair Maintenance. Discrete \& Computational Geometry, 19(2):175–195, 1998.
\bibitem{r14} Harry Buhrman, Christoph Durr, Mark Heiligman, Peter Hoyer, Frédéric Magniez, Miklos Santha, and Ronald De Wolf. Quantum Algorithms for Element Distinctness. In Proceedings 16th Annual IEEE Conference on Computational Complexity, pages 131–137. IEEE, 2001.
\bibitem{r15} Harry Buhrman, Subhasree Patro, and Florian Speelman. The Quantum Strong ExponentialTime Hypothesis. arXiv preprint, 2019. arXiv:1911.05686.
\bibitem{r16} Timothy M Chan and Ryan Williams. Deterministic APSP, orthogonal vectors, and more: Quickly derandomizing Razborov-Smolensky. In Proceedings of the twenty-seventh annual ACM-SIAM symposium on Discrete algorithms, pages 1246–1255. Society for Industrial and Applied Mathematics, 2016.
\bibitem{r17} Lijie Chen. On the Hardness of Approximate and Exact (Bichromatic) Maximum Inner Product. arXiv preprint, 2018. arXiv:1802.02325.
\bibitem{r18} Lijie Chen and Ryan Williams. An equivalence class for orthogonal vectors. In Proceedings of the Thirtieth Annual ACM-SIAM Symposium on Discrete Algorithms, pages 21–40. SIAM, 2019.
\bibitem{r19} Kenneth L Clarkson. A Randomized Algorithm for Closest-Point Queries. SIAM Journal on Computing, 17(4):830–847, 1988.
\bibitem{r20} Thomas H Cormen, Charles E Leiserson, Ronald L Rivest, and Clifford Stein. Introduction to Algorithms. MIT press, 2009.
\bibitem{r21} Evgeny Dantsin, Vladik Kreinovich, and Alexander Wolpert. On Quantum Versions of Record-breaking Algorithms for SAT. SIGACT News, 36(4):103–108, December 2005. doi: 10.1145/1107523.1107524.
\bibitem{r22} R. David, K. S., and B. Laekhanukit. On the Complexity of Closest Pair via Polar-Pair of Point-Sets. SIAM Journal on Discrete Mathematics, 33(1):509–527, 2019.
\bibitem{r23} Christoph Durr and Peter Hoyer. A quantum algorithm for finding the minimum. arXiv preprint, 1996. arXiv:quant-ph/9607014.
\bibitem{r24} Lov K. Grover. A Fast Quantum Mechanical Algorithm for Database Search. In Proceedings of the Twenty-eighth ACM Symposium on Theory of Computing - STOC '96. ACM Press, 1996. doi:10.1145/237814.237866.
\bibitem{r25} Timon Hertli. Improved Exponential Algorithms for SAT and ClSP. PhD thesis, ETH Zurich, 2015.
\bibitem{r26} Russell Impagliazzo and Ramamohan Paturi. On the Complexity of k-SAT. J. Comput. Syst. Sci., 62(2):367–375, March 2001. doi:10.1006/jcss.2000.1727.
\bibitem{r27} Toshiya Itoh, Tatsuya Nagatani, and Jun Tarui. Explicit Construction for k-Wise Nearly Random Permutations by Iterated Feistel Transform. Workshop on Randomness and Computation, 2005.
\bibitem{r28} Stacey Jeffery. Frameworks for Quantum Algorithms. PhD thesis, University of Waterloo, 2014.
\bibitem{r29} CS Karthik and Pasin Manurangsi. On Closest Pair in Euclidean Metric: Monochromatic is as Hard as Bichromatic. 10th Innovations in Theoretical Computer Science, 2019.
\bibitem{r30} S. Khuller and Y. Matias. A Simple Randomized Sieve Algorithm for the Closest-Pair Problem. Information and Computation, 118(1):34–37, 1995.
\bibitem{r31} Victor Klee. On the Complexity of d-Dimensional Voronoi Diagrams. Archiv der Mathematik, 34(1):75–80, 1980. doi:10.1007/bf01224932.
\bibitem{r32} Jon Kleinberg and Eva Tardos. Algorithm Design. Pearson Education India, 2006.
\bibitem{r33} Frédéric Magniez, Ashwin Nayak, Jérémie Roland, and Miklos Santha. Search via Quantum Walk. SIAM Journal on Computing, 40(1):142–164, 2011.
\bibitem{r34} Chris Marriott and John Watrous. Quantum Arthur–Merlin games. Computational Complexity, 14(2):122–152, 2005. doi:10.1007/s00037-005-0194-x.
\bibitem{r35} Abdullah Mueen, Eamonn Keogh, Qiang Zhu, Sydney Cash, and Brandon Westover. Exact Discovery of Time Series Motifs. In Proceedings of the 2009 SIAM international conference on data mining, pages 473–484. SIAM, 2009.
\bibitem{r36} Alexandros Nanopoulos, Yannis Theodoridis, and Yannis Manolopoulos. C2P: Clustering based on Closest Pairs. In VLDB, 2001.
\bibitem{r37} Michael A Nielsen and Isaac Chuang. Quantum Computation and Quantum Information, 2002.
\bibitem{r38} Ramamohan Paturi and Pavel Pudlák. On the Complexity of Circuit Satisfiability. In Proceedings of the forty-second ACM symposium on Theory of computing, pages 241–250. ACM, 2010.
\bibitem{r39} Ramamohan Paturi, Pavel Pudlák, Michael E Saks, and Francis Zane. An Improved Exponential-Time Algorithm for k-SAT. Journal of the ACM (JACM), 52(3):337–364, 2005.
\bibitem{r40} Michael O Rabin. Probabilistic Algorithms Algorithms and Complexity: New Directions and Recent Results, 1976.
\bibitem{r41} Kunihiko Sadakane, Norito Sugawara, and Takeshi Tokuyama. Quantum Algorithms for Intersection and Proximity Problems. In Peter Eades and Tadao Takaoka, editors, Algorithms and Computation, pages 148–159, Berlin, Heidelberg, 2001. Springer Berlin Heidelberg.
\bibitem{r42} Dominik Scheder and John P Steinberger. PPSZ for General k-SAT-Making Hertli’s Analysis Simpler and 3-SAT Faster. In 32nd Computational Complexity Conference (CCC 2017). Schloss Dagstuhl-Leibniz-Zentrum fuer Informatik, 2017.
\bibitem{r43} T Schoning. A Probabilistic Algorithm for k-SAT and Constraint Satisfaction Problems. In 40th Annual Symposium on Foundations of Computer Science (Cat. No. 99CB37039), pages 410–414. IEEE, 1999.
\bibitem{r44} M. I. Shamos and D. Hoey. Closest-Point Problems. In 16th Annual Symposium on Foundations of Computer Science (sfcs 1975), pages 151–162, 1975.
\bibitem{r45} Mario Szegedy. Quantum Speed-Up of Markov Chain Based Algorithms. In 45th Annual IEEE symposium on foundations of computer science, pages 32–41. IEEE, 2004.
\bibitem{r46} Virginia Vassilevska Williams. Hardness of Easy Problems: Basing Hardness on Popular Conjectures such as the Strong Exponential Time Hypothesis (invited talk). In 10th International Symposium on Parameterized and Exact Computation (IPEC 2015). Schloss Dagstuhl-LeibnizZentrum fuer Informatik, 2015.
\bibitem{r47} Nilton Volpato and Arnaldo Moura. A fast quantum algorithm for the closest bichromatic pair problem, 2010.
\bibitem{r48} Ryan Williams. Pairwise comparison of bit vectors. Theoretical Computer Science Stack Exchange. URL: https://cstheory.stackexchange.com/q/37369.
\bibitem{r49} Ryan Williams. A New Algorithm for Optimal 2-Constraint Satisfaction and Its Implications. Theoretical Computer Science, 348(2-3):357–365, 2005.
\bibitem{r50} Ryan Williams and Huacheng Yu. Finding Orthogonal Vectors In Discrete Structures. In Proceedings of the twenty-fifth annual ACM-SIAM symposium on Discrete algorithms, pages 1867–1877. SIAM, 2014.
\bibitem{r51} Raymond Chi-Wing Wong, Yufei Tao, Ada Wai-Chee Fu, and Xiaokui Xiao. On Efficient Spatial Matching. In Proceedings of the 33rd International Conference on Very Large Data Bases, VLDB ’07, pages 579–590, 2007.
\bibitem{r52} A. C. . Yao. Lower Bounds for Algebraic Computation Trees with Integer Inputs. In 30th Annual Symposium on Foundations of Computer Science, pages 308–313, 1989.
\end{thebibliography}
\end{document}
